% !TEX program = uplatex
\documentclass[uplatex,dvipdfmx,a4paper,10.5pt]{jsarticle}

\usepackage[margin=20mm]{geometry}
\usepackage{amsmath,amssymb,amsthm,mathtools}
\usepackage{mathpartir}
\usepackage{booktabs,array,tabularx}
\usepackage{enumitem}
\usepackage{xcolor}
\usepackage[most]{tcolorbox}
\usepackage{hyperref}
\usepackage{pxjahyper}
\hypersetup{
  colorlinks=true,
  linkcolor=blue!45!black,
  citecolor=blue!45!black,
  urlcolor=blue!45!black,
  pdftitle={Constructor-free Directed Weight Solver and Termination},
  pdfsubject={Row-only weighted bounds graph, residual descent, and cycle-neutral termination},
  pdfauthor={}
}

\setlength{\parindent}{1em}
\setlength{\parskip}{0.18em}
\setlist[itemize]{topsep=0.35em,itemsep=0.14em,parsep=0pt}
\setlist[enumerate]{topsep=0.35em,itemsep=0.14em,parsep=0pt}
\allowdisplaybreaks

\definecolor{deepblue}{RGB}{36,70,115}
\definecolor{softblue}{RGB}{238,245,252}
\definecolor{softgreen}{RGB}{239,248,241}
\definecolor{softyellow}{RGB}{255,249,226}
\definecolor{softred}{RGB}{253,239,239}
\newtcolorbox{formalbox}[1][]{
  breakable,colback=softblue,colframe=deepblue!70,boxrule=0.55pt,
  arc=1.2mm,title={形式化},fonttitle=\bfseries,#1}
\newtcolorbox{notebox}[1][]{
  breakable,colback=softyellow,colframe=orange!70!black,boxrule=0.55pt,
  arc=1.2mm,title={注意},fonttitle=\bfseries,#1}
\newtcolorbox{statusbox}[1][]{
  breakable,colback=softred,colframe=red!55!black,boxrule=0.55pt,
  arc=1.2mm,title={結論},fonttitle=\bfseries,#1}
\newtcolorbox{intuitionbox}[1][]{
  breakable,colback=softgreen,colframe=green!45!black,boxrule=0.55pt,
  arc=1.2mm,title={見方},fonttitle=\bfseries,#1}

\newtheoremstyle{jpplain}
  {0.65em}{0.65em}{\normalfont}{}
  {\bfseries}{。}{0.5em}{}
\theoremstyle{jpplain}
\newtheorem{definition}{定義}[section]
\newtheorem{lemma}[definition]{補題}
\newtheorem{proposition}[definition]{命題}
\newtheorem{theorem}[definition]{定理}
\newtheorem{corollary}[definition]{系}
\newtheorem{invariant}[definition]{不変条件}
\newtheorem{example}[definition]{例}
\newtheorem{remark}[definition]{注意}
\renewcommand{\proofname}{証明}

% symbols
\newcommand{\Caps}{\mathcal B}
\newcommand{\Id}{\mathsf{Id}}
\newcommand{\Var}{\mathsf{Var}}
\newcommand{\Nat}{\mathbb N}
\newcommand{\topc}{\top_{\Caps}}
\newcommand{\botc}{\bot_{\Caps}}
\newcommand{\meet}{\mathbin{\wedge}}
\newcommand{\join}{\mathbin{\vee}}
\newcommand{\bigmeet}{\mathop{\bigwedge}}
\newcommand{\diff}{\mathbin{\setminus}}
\newcommand{\Weight}{\mathcal W}
\newcommand{\Pair}{\mathcal P}
\newcommand{\idw}{\mathbf 1_{\Weight}}
\newcommand{\idp}{\mathbf 1_{\Pair}}
\newcommand{\wcomp}{\mathbin{\cdot}}
\newcommand{\pcomp}{\mathbin{\diamond}}
\newcommand{\flatten}{\mathsf{flat}}
\newcommand{\amb}{\mathsf{amb}}
\newcommand{\armed}{\mathsf{armed}}
\newcommand{\budget}{\mathsf{budget}}
\newcommand{\common}{\mathsf{common}}
\newcommand{\res}{\mathsf{res}}
\newcommand{\respair}{\widehat{\mathsf{res}}}
\newcommand{\seed}{\chi}
\newcommand{\capheight}{\mathsf{ht}}
\newcommand{\nf}{\mathsf{nf}}
\newcommand{\seen}{\mathsf{Seen}}
\newcommand{\memo}{\mathsf{Memo}}
\newcommand{\row}[2]{[#1;#2]}
\newcommand{\edge}[3]{#1\xrightarrow{#2}#3}
\newcommand{\epedge}[4]{#1\xrightarrow{#2}[#3;#4]}
\newcommand{\rid}{\mathsf{resid}}
\newcommand{\supp}{\mathsf{supp}}

\title{型コンストラクタを除いた方向付き weight solver の停止性\\
\large row-only bounds graph、residual budget、cycle-neutral saturation}
\author{}
\date{2026年6月21日}

\begin{document}
\maketitle

\begin{abstract}
本稿は、方向付き stack weight の停止性だけを切り出すため、値型、関数型、ADT、computation 型、
constructor 分解、項の制約生成、compact extraction を全て外した row-only core を定義する。
型 constructor の情報は、有限個の row variable と row endpoint からなる初期 graph へ既に atomize されたものとする。

各 subtract id には固定 seed を与え、filter-free な一 id の作用を
未対応 pop 数、残る push 数、永続 floor の三つ組で表す。
さらに ordinary row head の二重消費も同じ下降測度で扱うため、weight 全体に ambient residual budget を一つ持たせる。
左右 label は $\Weight\times\Weight^{\mathrm{op}}$ として合成し、row endpoint で一つの weight へ flatten する。

exact worklist の無条件停止性は偽である。非単位 label を持つ自己 cycle は
$b,b^2,b^3,\ldots$ を生成する。
そこで、各 directed cycle の label 積が単位元である cycle-neutrality を明示的に要求する。
有限 capability algebra、cycle-neutrality、exact cache、residual memoization の下で、
path label は有限個、nonempty residual step は capability lattice の高さ以下、生成 node と seen key は有限個となる。
従って worklist solver は停止する。
\end{abstract}

\tableofcontents

\section{簡略化の範囲}

この稿では次を意図的に記述しない。

\begin{itemize}
  \item 値型、関数型、ADT、computation 型の grammar。
  \item W-Fun、W-Comp、constructor variance の分解規則。
  \item 表面項、let 多相、\texttt{let rec}、compact Pos/Neg 表示。
  \item colored operational semantics。
\end{itemize}

残すのは、有限な row variable graph、方向付き edge label、variable replay、weighted row split だけである。
constructor 分解は別の有限 pre-pass とし、その出力を本稿の初期 graph とする。
反変位置での左右交換も pre-pass 側で済ませ、core は既に得られた pair label を扱う。

\begin{statusbox}
この簡略化の下でも、停止性は無条件には成り立たない。
exact natural-number counter を保つ限り、非中立 cycle を拒否するか、別定理を持つ widening を導入する必要がある。
本文では exact solver を保ち、cycle-neutrality を検査して非中立 cycle を拒否する。
\end{statusbox}

\section{有限 capability algebra と weight}

\subsection{capability algebra}

\begin{definition}[有限 capability algebra]
$\Caps$ を有限 Boolean algebra とする。
単位元を $\topc$、零元を $\botc$、meet、join、difference をそれぞれ
$\meet,\join,\diff$ と書く。
その高さを $\capheight(\Caps)$ とする。
\end{definition}

集合モデルでは $\Caps=\mathcal P(\mathsf{Fam})$ であり、
$A\diff U=A\cap\overline U$ である。
本文の停止性証明が使うのは、有限性と
\[
 (A\diff U)\meet(B\diff U)=(A\meet B)\diff U
\]
だけである。

\subsection{一 id の正規形}

subtract id の有限集合を $\Id$ とし、各 id に固定 seed
\[
  \seed:\Id\to\Caps
\]
を与える。

\begin{definition}[entry]
一 id $u$ の filter-free summary を
\[
  e=(p,q,F)\in\Nat\times\Nat\times\Caps
\]
とする。
$p$ は入力側へ要求する未対応 pop 数、$q$ は最後に残る local push 数、
$F$ はその id に永続する residual floor である。
単位 entry は $(0,0,\topc)$ である。
\end{definition}

active push が存在するとき、その全 frame は同じ representative capability
$\seed(u)\meet F$ を持つ。
従って family の列は不要で、個数 $q$ と floor $F$ だけで足りる。

\begin{definition}[entry composition]
$e=(p,q,F)$、$f=(r,s,G)$ とし、$c=\min(q,r)$ と置く。
左から $e$、次に $f$ を作用させる積を
\[
  e\wcomp f
  =
  (p+r-c,\ q+s-c,\ F\meet G)
\]
と定める。
\end{definition}

\begin{lemma}[entry monoid]
entry は上の積と $(0,0,\topc)$ により monoid をなす。
\end{lemma}

\begin{proof}
push/pop 部分は関係 $\mathsf{push}\,\mathsf{pop}=1$ を持つ bicyclic monoid の標準形
$\mathsf{pop}^{p}\mathsf{push}^{q}$ である。
$c=\min(q,r)$ は二語の境界で相殺される個数に一致する。
floor 成分は meet monoid であり、二成分の直積なので結合律と単位律を得る。
\end{proof}

\subsection{weight と ambient residual budget}

ordinary row 自身が既に消費した family を、active id が無い場合にも忘れないため、
weight に ambient budget を一つ持たせる。
これは runtime color ではなく、row residual の静的な残量である。

\begin{definition}[weight]
weight は
\[
  W=(A,(e_u)_{u\in\Id})
\]
である。$A\in\Caps$ を ambient budget とし、$e_u$ は有限台 entry map とする。
積は
\[
  (A,(e_u))\wcomp(B,(f_u))
  =
  (A\meet B,(e_u\wcomp f_u))
\]
である。単位 weight は
\[
  \idw=(\topc,(0,0,\topc)_{u\in\Id})
\]
とする。
\end{definition}

\begin{lemma}[weight monoid]
$(\Weight,\wcomp,\idw)$ は monoid である。
また
\[
  \amb(W\wcomp V)=\amb(W)\meet\amb(V)\leq\amb(W)
\]
である。
\end{lemma}

\begin{proof}
ambient は meet monoid、各 id は entry monoid なので pointwise に従う。
後半は meet の単調性による。
\end{proof}

\subsection{左右 label を一個の pair として扱う}

\begin{definition}[pair weight]
\[
  \Pair=\Weight\times\Weight^{\mathrm{op}}
\]
とし、要素を $b=\langle L\mid R\rangle$ と書く。
積と単位元を
\[
 \langle L\mid R\rangle\pcomp\langle L'\mid R'\rangle
 =\langle L\wcomp L'\mid R'\wcomp R\rangle,
 \qquad
 \idp=\langle\idw\mid\idw\rangle
\]
で定める。
row endpoint で観測する weight は
\[
  \flatten(\langle L\mid R\rangle)=L\wcomp R
\]
である。
\end{definition}

\begin{lemma}[pair monoid]
$(\Pair,\pcomp,\idp)$ は monoid である。
\end{lemma}

\begin{proof}
左成分は $\Weight$、右成分は opposite monoid の積である。
\end{proof}

\section{common budget と residual の strict descent}

$W=(A,(p_u,q_u,F_u)_u)$ とする。

\begin{definition}[armed id と entry budget]
\[
 \armed(W)
 =\{u\mid p_u>0\text{ or }q_u>0\text{ or }F_u\neq\topc\}.
\]
\[
 B_u(W)=
 \begin{cases}
   \seed(u)\meet F_u,&q_u>0,\\
   F_u,&q_u=0.
 \end{cases}
\]
\end{definition}

pop-only entry は $B_u=\topc$ を寄与するので、subtraction budget を狭めない。

\begin{definition}[common budget]
\[
  \common(W)
  =\amb(W)\meet
   \bigmeet_{u\in\armed(W)}B_u(W).
\]
空 meet は $\topc$ とする。
\end{definition}

ambient を常に meet へ含めるため、active id が無くても
$\common(W)=\amb(W)$ となる。

\begin{definition}[residual]
$\botc\neq U\leq\common(W)$ とする。
residual floor weight $F_U(W)$ を次で定める。

\begin{itemize}
  \item ambient 成分は $\amb(W)\diff U$。
  \item $u\in\armed(W)$ の entry は $(0,0,B_u(W)\diff U)$。
  \item それ以外の entry は単位 entry。
\end{itemize}

\[
  \res_U(W)=W\wcomp F_U(W).
\]
\end{definition}

\begin{lemma}[一 entry の residual]
$u\in\armed(W)$ とする。
$\res_U(W)$ における $u$ の budget は
\[
  B_u(W)\diff U
\]
であり、未対応 pop 数と active push 数は変わらない。
\end{lemma}

\begin{proof}
$q_u>0$ なら新 floor は
\[
 F_u\meet(B_u\diff U)=B_u\diff U
\]
であり、$B_u\diff U\leq\seed(u)$ なので新 representative は
$B_u\diff U$ である。
$q_u=0$ なら $B_u=F_u$ なので同じ計算になる。
residual entry の pop/push はともに 0 である。
\end{proof}

\begin{theorem}[budget conservation]
\label{thm:budget-conservation}
$\botc\neq U\leq\common(W)$ なら
\[
  \amb(\res_U(W))=\amb(W)\diff U,
\]
\[
  \common(\res_U(W))=\common(W)\diff U.
\]
\end{theorem}

\begin{proof}
ambient の式は定義から直接従う。
前補題により各 armed entry budget は $B_u\diff U$ になる。
従って Boolean algebra の分配則から
\[
\begin{split}
 \common(\res_U(W))
 &= (\amb(W)\diff U)
    \meet\bigmeet_u(B_u(W)\diff U)\\
 &= \left(\amb(W)\meet\bigmeet_uB_u(W)\right)\diff U\\
 &= \common(W)\diff U.
\end{split}
\]
\end{proof}

\begin{corollary}[strict residual descent]
\label{cor:strict-descent}
$\botc\neq U\leq\common(W)$ なら
\[
  \amb(\res_U(W))<\amb(W).
\]
従って、一つの provenance path に現れる nonempty residual step の個数は
$\capheight(\Caps)$ 以下である。
\end{corollary}

\begin{proof}
$U\leq\common(W)\leq\amb(W)$ かつ $U\neq\botc$ なので、
$\amb(W)\diff U$ は $\amb(W)$ の真の下位元である。
weight をさらに右へ合成しても ambient は meet により増えない。
従って residual step ごとに strict descending chain を得る。
\end{proof}

\begin{corollary}[二重消費禁止]
$U=H\meet\common(W)\neq\botc$ なら
\[
  H\meet\common(\res_U(W))=\botc.
\]
\end{corollary}

\begin{proof}
Theorem~\ref{thm:budget-conservation} から
\[
 H\meet(\common(W)\diff U)
 =(H\meet\common(W))\diff U
 =U\diff U=\botc.
\]
\end{proof}

\section{constructor-free bounds graph machine}

\subsection{primitive fact と control graph}

初期 row variable 集合 $V_0\subseteq\Var$ と初期 fact 集合は有限とする。
型 constructor の分解は既に終わっており、primitive fact は次の二形だけである。

\[
  \edge{x}{b}{y},
  \qquad
  \epedge{x}{b}{H}{y}
  \qquad
  (x,y\in V_0,\ b\in\Pair,\ H\in\Caps).
\]

第一は primitive variable edge、第二は primitive row endpoint である。
Replay が追加する推移 edge と伝播 endpoint は \emph{derived fact} と呼び、primitive fact table とは分けて保存する。

\begin{definition}[control graph]
primitive fact 集合 $P$ の control graph $\mathsf{Ctl}(P)$ を次で定める。

\begin{itemize}
  \item primitive variable edge $\edge{x}{b}{y}$ ごとに arc $\edge{x}{b}{y}$ を置く。
  \item primitive endpoint $\epedge{x}{b}{H}{y}$ ごとに、head $H$ を忘れた tail projection arc
  $\edge{x}{b}{y}$ を置く。
\end{itemize}
\end{definition}

endpoint の tail projection を最初から control graph に入れるのは、$U=\botc$ の split が後で同じ label の
variable edge を作る可能性を先に数えるためである。これは実行可能 path の過大近似だが、停止性には安全である。

worklist state は queue、exact seen set、primitive fact table、derived fact table、residual memo table からなる。
seen key は source、target、head、正規化済み pair label のうち、その fact に現れる全成分を含む。

\subsection{処理規則}

\paragraph{Replay.}
\[
  \edge{x}{b}{y},\quad \edge{y}{c}{z}
  \quad\Longrightarrow\quad
  \edge{x}{b\pcomp c}{z},
\]
\[
  \edge{x}{b}{y},\quad \epedge{y}{c}{H}{z}
  \quad\Longrightarrow\quad
  \epedge{x}{b\pcomp c}{H}{z}.
\]
生成物は derived fact table へ置く。
従って任意の derived variable edge は control graph の path label であり、任意の derived endpoint は
control path と primitive endpoint の積である。これは lower/upper bound replay を path composition として
書いたものに等しい。

\paragraph{Weighted split.}
primitive または derived endpoint $\epedge{x}{b}{H}{y}$ に対し
$W=\flatten(b)$、
\[
  U=H\meet\common(W)
\]
と置く。

\begin{enumerate}
  \item $U=\botc$ なら $\edge{x}{b}{y}$ を derived fact として enqueue する。
  この edge の arc は endpoint の tail projection として既に control graph に含まれている。
  \item $U\neq\botc$ なら
  \[
    \respair_U(b)=\langle\res_U(\flatten(b))\mid\idw\rangle
  \]
  とし、canonical residual node
  \[
    \gamma=\rid(x,U,\respair_U(b))
  \]
  を memo table から得る。新しい key なら primitive row head endpoint
  \[
    \epedge{x}{\idp}{U}{\gamma}
  \]
  を一度だけ挿入する。key の新旧にかかわらず、この endpoint の target $y$ に対する primitive residual continuation edge
  \[
    \edge{\gamma}{\respair_U(b)}{y}
  \]
  が未登録なら挿入し、未処理の fact を enqueue する。
\end{enumerate}

同じ $(x,U,\respair_U(b))$ には同じ $\gamma$ を返す。target $y$ は memo key に含めない。
row endpoint を越えた後は本稿に constructor/variance が無いため、residual weight を左成分へ正規化してよい。

\begin{invariant}[residual provenance]
nonempty split が生成する primitive continuation edge の label は、split 前の endpoint label より
ambient budget が真に小さい。その edge より後ろへ control path を合成しても、この減少は復活しない。
\end{invariant}

\begin{proof}
continuation label は $\respair_U(b)$ であり、Corollary~\ref{cor:strict-descent} により
\[
  \amb(\res_U(\flatten(b)))<\amb(\flatten(b)).
\]
control path の後続 label を右へ合成しても、ambient 成分は meet されるだけなので増えない。
\end{proof}

\section{無条件停止性が偽である理由}

\begin{example}[非中立自己 cycle]
一変数 $x$ を取り、$b$ を「ある id に一個の unmatched pop を持ち、他成分は単位元」である pair とする。
このとき $b\neq\idp$ であり、primitive edge
\[
  \edge{x}{b}{x}
\]
を入れると、Replay は
\[
  \edge{x}{b^2}{x},\quad
  \edge{x}{b^3}{x},\quad\ldots
\]
を順に生成する。$b^n$ の正規形の pop count は $n$ なので全て異なる。
exact seen set はこれらを別 key として扱うため、worklist は停止しない。
\end{example}

この反例は型 constructor と無関係である。
従って constructor の記述を外すだけでは停止性は得られず、cycle label に条件が必要になる。

\section{cycle-neutrality と path label の有限性}

\begin{definition}[cycle-neutral fact set]
primitive fact 集合 $P$ の control graph $\mathsf{Ctl}(P)$ の任意の directed cycle について、
arc label を順に積んだ結果が $\idp$ であるとき、$P$ を cycle-neutral と呼ぶ。
\end{definition}

実装は nonempty split が二 primitive fact を追加するたび control graph の SCC を更新し、
新しい cycle が neutral か検査する。
より強い十分条件として「SCC 内の非単位 arc を禁止」を採用してもよい。
非 neutral cycle は counter を飽和して続けず、\texttt{UnbalancedEffectBoundary} として拒否する。

\begin{lemma}[closed walk の中立性]
有限 cycle-neutral control graph の任意の closed walk の label は $\idp$ である。
\end{lemma}

\begin{proof}
closed walk が simple cycle なら定義による。
そうでなければ途中で同じ頂点を再訪するので、その間のより短い closed subwalk を取り出せる。
長さに関する帰納法で subwalk の label は単位元であり、それを除いた closed walk も単位元である。
\end{proof}

\begin{lemma}[control path label の有限性]
\label{lem:path-finite}
有限 cycle-neutral control graph では、固定した二頂点間の path に現れる正規化済み pair label は有限個である。
\end{lemma}

\begin{proof}
頂点を二度通る path から、その間の closed subwalk を除く。
前補題により除いた部分の label は $\idp$ なので、全体 label は変わらない。
これを繰り返すと simple path になる。有限 graph の simple path は有限個である。
\end{proof}

\begin{notebox}
cycle-neutrality は必要条件ではなく、exact termination を簡潔に保証する十分条件である。
例えば冪等な floor-only cycle は非単位でも生成 label 集合が有限な場合がある。
より一般には「各 SCC の closed-walk label が生成する submonoid が有限」を仮定すれば同じ証明が通る。
本文は実装で判定しやすい cycle-neutrality に固定する。
\end{notebox}

\section{生成状態の有限性}

\subsection{residual phase normal form}

control path だけを進む区間を \emph{path phase}、nonempty weighted split が residual node を一個生成する遷移を
\emph{residual phase} と呼ぶ。任意の生成履歴は概念的に
\[
  \lambda_0;\res_{U_1};\lambda_1;\res_{U_2};\cdots;
  \res_{U_k};\lambda_k
\]
と分解できる。ここで各 $\lambda_j$ は residual step を含まない control path の label である。

\begin{lemma}[residual phase 数の上界]
任意の生成履歴で $k\leq\capheight(\Caps)$ である。
\end{lemma}

\begin{proof}
各 residual phase は $U_j\neq\botc$ の場合にだけ生成される。
Corollary~\ref{cor:strict-descent} と residual provenance invariant により、履歴上の ambient budget は
residual phase ごとに真に減り、path phase では増えない。
有限 lattice の strict descending chain の長さは高さ以下である。
\end{proof}

\begin{lemma}[深さごとの有限生成]
\label{lem:finite-by-depth}
residual phase を高々 $k$ 個含む primitive node、primitive fact、derived fact、seen key の集合は有限である。
\end{lemma}

\begin{proof}
$k$ に関する帰納法。

$k=0$ では primitive node と primitive fact は有限な初期集合だけである。
その control graph は有限かつ cycle-neutral なので、Lemma~\ref{lem:path-finite} により control path label は有限個である。
各 derived variable edge はそのいずれか、各 derived endpoint は有限個の path label と primitive endpoint の積なので有限個である。

$k$ 以下で primitive node と primitive fact が有限と仮定する。
その有限 cycle-neutral control graph 上の path label は再び有限個で、従って derived fact と endpoint key も有限個である。
一つの endpoint key から split が作る $U$ は有限集合 $\Caps$ の元なので、
residual memo key $(x,U,\respair_U(b))$ の候補は有限個である。
同じ key は一 node に hash-cons されるため、深さ $k+1$ の新 node は有限個である。
各新 node は primitive row head endpoint を一つだけ持ち、continuation edge は有限個の target ごとに高々一つである。
従って深さ $k+1$ までの primitive fact 集合も有限であり、cycle-neutrality の仮定の下で次の帰納段階へ進める。
exact seen key は有限な fact 成分の直積なので有限である。
\end{proof}

\begin{theorem}[生成状態の有限性]
\label{thm:state-finite}
有限 capability algebra、有限初期 primitive fact 集合、各段階の cycle-neutrality、residual memoization の下で、
solver が生成しうる primitive node、primitive fact、derived fact、row endpoint、seen key は全て有限個である。
\end{theorem}

\begin{proof}
residual phase 数は前補題により $\capheight(\Caps)$ 以下である。
Lemma~\ref{lem:finite-by-depth} を $k=\capheight(\Caps)$ まで適用すれば結論を得る。
\end{proof}

\section{worklist solver の停止性}

\begin{theorem}[constructor-free solver termination]
\label{thm:termination}
次を仮定する。

\begin{enumerate}
  \item $\Caps$、$\Id$、初期 row variable、初期 primitive edge、初期 primitive endpoint は有限である。
  \item 全 weight と pair は処理前に一意な正規形へ canonicalize される。
  \item seen key は fact の全成分と正規化済み pair label を含み、同じ key は一度だけ処理する。
  \item 初期 control graph は cycle-neutral で、nonempty split による primitive fact 追加後も cycle-neutrality を保つ。
  \item residual node は $(x,U,\respair_U(b))$ で memoize する。
\end{enumerate}

このとき constructor-free worklist solver は有限回で停止する。
\end{theorem}

\begin{proof}
Theorem~\ref{thm:state-finite} により、queue へ入りうる exact key の全体は有限集合である。
各 key は seen set により高々一回だけ本体処理される。
一回の処理が列挙する既存 fact は有限で、一つの key から追加される候補も有限個である。
従って本体処理の総回数は有限で、queue は最終的に空になる。
\end{proof}

\begin{corollary}[飽和点の一意性]
処理順序に依存せず、solver の停止状態は初期制約を含み Replay と weighted split で閉じた state のうち最小である。
\end{corollary}

\begin{proof}
各規則は既存 fact から一意に決まる fact だけを追加する単調な closure operator である。
有限状態空間上で fair worklist が全規則を処理するため、得られる fixed point は最小閉包である。
\end{proof}


\section{元の体系へ戻すときの境界}

\begin{intuitionbox}
この停止性定理が直接扱うのは、constructor 分解後の row-only graph である。
関数、ADT、computation 型を戻す場合は、それらの分解が有限個の hash-consed subnode しか作らないことと、
分解後の初期 primitive fact が作る control graph が cycle-neutral であることを別補題として示せばよい。
constructor 自体は停止測度に入らず、問題になるのは variable cycle の label 増幅と residual の無限生成だけである。
\end{intuitionbox}

実装上の義務は次の五点に集約される。

\begin{enumerate}
  \item cache 前に entry、weight、pair を必ず canonicalize する。
  \item right component は opposite order で replay する。
  \item primitive fact の追加で non-neutral SCC が生じたら exact core では拒否する。
  \item residual key を source、nonempty $U$、residual pair で共有し、target tail は key に入れない。
  \item ambient residual budget を runtime marker に materialize しない。これは ordinary row の消費済み集合を記録する静的成分である。
\end{enumerate}

\begin{statusbox}
型 constructor の記述を外した core では、停止性の未解決部分は完全に二つへ分離できる。

\begin{enumerate}
  \item variable replay: cycle-neutrality により path label を有限化する。
  \item row residual: ambient budget の strict descent により生成深さを有限化する。
\end{enumerate}

この二つと exact cache、residual memoization から worklist termination が従う。
非 neutral cycle を許したまま exact counter を保つ無条件停止性は、自己 cycle の反例により成り立たない。
\end{statusbox}

\begin{thebibliography}{9}
\bibitem{parreaux2020}
Lionel Parreaux.
\newblock The Simple Essence of Algebraic Subtyping: Principal Type Inference with Subtyping Made Easy.
\newblock PACMPL 4(ICFP), 2020.

\bibitem{damasmilner1982}
Luis Damas and Robin Milner.
\newblock Principal type-schemes for functional programs.
\newblock POPL, 1982.
\end{thebibliography}

\end{document}
